\chapter{Introduction}

Software quality assurance has automated a great deal of what was once
manual regression testing, yet one activity has resisted that
automation: the authoring of the end-to-end test scripts themselves.
Turning an informal product requirement into a syntactically precise,
assertion-rich automation script --- one that exercises both the
happy-path journey a user is expected to take through a feature and the
edge cases most likely to expose a defect {[}11{]} --- remains a task
performed by a human engineer, and one that is repetitive,
time-consuming, and sensitive to that engineer's familiarity with the
specific application domain and test framework in use. It is this gap
between a plainly stated requirement and a working test script that
motivates the present research.

Large language models trained on code have shown that a pre-trained
transformer can produce output that is structurally plausible --- code
that reads as syntactically correct even where it is not necessarily
logically correct --- from a natural-language description alone
{[}12{]}. A general-purpose model prompted with no task-specific
grounding, however, has no exposure to the particular selector
conventions, assertion idioms, and API surface mandated by a specific
test framework such as Cypress {[}13{]} or Playwright {[}14{]};
zero-shot or few-shot prompting {[}12{]} of such a model typically
yields scripts that are structurally incomplete, rely on deprecated
APIs, or omit the assertions needed for the test to carry any real
verification value. Closing this gap requires more than a better prompt:
it requires the model to be grounded in the target framework's
conventions and given a mechanism to check and correct its own output.

This dissertation makes three complementary contributions toward that
goal. A domain-specific dataset of 279 aligned
user-story-and-test-script pairs was constructed through a structured,
quality-validated generation pipeline, split evenly and independently
between the Cypress and Playwright frameworks. Two open-weight language
models, Phi-3 Mini and Gemma 4 E4B, were then adapted to this dataset
under the QLoRA (Quantised Low-Rank Adaptation) parameter-efficient
fine-tuning paradigm, producing one specialised adapter per model per
framework so that no adapter is ever exposed to the syntax of the
framework it was not trained for. Finally, an agentic
Build--Measure--Assess--Decide (BMAD) correction loop was designed and
implemented to autonomously generate, statically inspect, score, and ---
where necessary --- iteratively revise each script until it meets a
defined quality threshold or a maximum number of correction attempts is
reached, removing the need for a human reviewer in the refinement cycle
itself.

The research sits at the intersection of AI-assisted software
engineering and test automation, with direct practical relevance to
organisations that maintain continuous integration pipelines under Agile
delivery practices, where the marginal cost of authoring and maintaining
end-to-end test coverage is a recurring constraint on release velocity.
All eight phases of the work --- dataset construction, baseline
evaluation, fine-tuning, agentic-loop implementation, and comparative
statistical evaluation --- have been completed at the time of this
report. Chapter 2 surveys the literature that motivates the specific
design choices made; Chapter 3 describes the system design and
architecture; Chapter 4 details the implementation; Chapter 5 reports
the complete comparative evaluation across all systems studied; and
Chapter 6 concludes with a synthesis of findings, limitations, and
directions for future work.

\chapter{Literature Review}

This chapter surveys three bodies of literature that together frame the
design of the system described in this dissertation: LLM-driven code and
test generation, parameter-efficient fine-tuning, and agentic
self-correcting AI architectures. Each is examined for what it
establishes and, more importantly, for what it leaves unaddressed with
respect to the specific problem this dissertation tackles.

The use of large language models for automated code generation has
progressed from general-purpose function synthesis toward increasingly
specialised, test-oriented applications. Chen et al. {[}1{]} established
the field's reference evaluation methodology with Codex and the
HumanEval benchmark, demonstrating that a decoder-only transformer
trained on large code corpora attains substantial competence at
function-level synthesis and popularising pass-at-k as the standard
accuracy metric for generated code. Austin et al. {[}2{]} broadened this
evaluation basis with MBPP, a benchmark of crowd-sourced programming
problems spanning a wider range of task types than HumanEval alone.
Closer to the concern of this dissertation, Liu et al. {[}3{]} examined
LLM-generated unit tests specifically, showing that a model supplied
with a method signature can produce tests that measurably improve code
coverage --- evidence that test-script generation, not only
application-code generation, is within reach of current models, provided
they are given adequate task grounding. StarCoder {[}4{]} and
DeepSeek-Coder {[}5{]} extended the open-weight code-model landscape by
training on broader multi-language corpora, yielding models that are
both competitive on standard benchmarks and, unlike proprietary
alternatives, available for further fine-tuning --- a property this
dissertation depends upon directly.

None of this literature, however, addresses UI test-script generation
for a named commercial framework with selector-level correctness as an
explicit criterion: HumanEval and MBPP evaluate general-purpose function
synthesis, and the unit-test generation work of Liu et al. targets
arbitrary application code rather than the domain-specific object model,
locator conventions, and assertion idioms of a browser-automation
framework such as Cypress or Playwright. Studies of hand-authored UI
test suites, such as Leotta et al.'s comparison of Selenium WebDriver
design patterns {[}15{]}, confirm that writing and maintaining such
scripts carries a substantial engineering cost in its own right ---
precisely the cost that automated generation seeks to reduce.

Because full-parameter fine-tuning of a billion-parameter model demands
GPU infrastructure well beyond consumer hardware, this dissertation
relies instead on parameter-efficient fine-tuning (PEFT), a family of
techniques that adapts a pre-trained model to a target task by updating
only a small fraction of its parameters while leaving the bulk of its
pre-trained knowledge frozen. Hu et al. {[}6{]} introduced Low-Rank
Adaptation (LoRA), which injects small trainable rank-decomposition
matrices alongside the frozen weights of each adapted layer, cutting the
number of trainable parameters by several orders of magnitude with
little measurable loss of task performance. Dettmers et al. {[}7{]}
extended this idea to QLoRA, quantising the frozen base model itself to
4-bit precision while training the LoRA adapters at higher precision, a
combination that for the first time made fine-tuning of
7-billion-parameter-class models practical on a single consumer GPU.
This is the specific technique this dissertation applies, and it is what
makes on-device fine-tuning on Apple Silicon --- rather than a cloud GPU
cluster --- a realistic option for a project of this scope.

Yet the QLoRA literature to date has been benchmarked and validated
almost exclusively on CUDA-capable NVIDIA hardware; Apple's Metal
Performance Shaders (MPS) backend, which this dissertation targets,
differs from CUDA in numerically significant ways --- for instance in
how mixed-precision gradient scaling behaves --- and its suitability for
producing a fine-tuned adapter of comparable quality has not been
previously reported for this task.

The third relevant body of work concerns agentic architectures, in which
a model does not merely produce a single output but reasons, acts, and
revises across an iterative loop, often incorporating feedback from an
external tool or from its own prior attempt. Yao et al. {[}8{]} proposed
ReAct, interleaving explicit reasoning traces with executable actions so
that a model can invoke external tools and incorporate their results
before producing a final answer. Wang et al. {[}9{]} demonstrated, with
AutoGen, that two cooperating agents --- one producing code and a second
critiquing it --- can improve code quality beyond what either agent
achieves alone. Shinn et al. {[}10{]} introduced Reflexion, in which an
agent inspects the outcome of a failed attempt and uses that reflection
to generate a corrected subsequent attempt without any change to the
underlying model weights. These three works collectively establish the
conceptual template followed by the BMAD loop developed in this
dissertation, in which an orchestrating process evaluates the syntactic
and semantic quality of a generated script and decides, on that basis,
whether to accept it or to trigger a further, feedback-guided generation
attempt.

None of these agentic frameworks, however, has been applied to the
specific problem of iteratively correcting a UI test script using
static-analysis feedback as the correction signal, nor evaluated with
the paired statistical rigour --- the same held-out items scored across
every system under comparison --- that this dissertation's evaluation
chapter applies.

Table 2.1 summarises the three gaps identified across this survey and
states how this dissertation addresses each of them; the corresponding
evidence for each closure is presented in Chapter 5.

\begin{longtable}[]{@{}lll@{}}
\caption{Literature gaps and how this dissertation addresses them}\tabularnewline
\toprule
\endhead
\textbf{Gap} & \textbf{Description} & \textbf{How Addressed} \\
Gap 1: No framework-idiomatic UI test-script fine-tuning & Prior code-
and test-generation work {[}1, 3{]} evaluates general-purpose function
or unit-test synthesis; none fine-tunes a model specifically on the
selector conventions and assertion idioms of a named commercial UI test
framework. & Two open-weight models were fine-tuned on a
framework-isolated, 279-pair Cypress/Playwright dataset, producing
framework-specific adapters evaluated against both zero-shot baselines
and independent syntax validation (Chapter 5). \\
Gap 2: QLoRA feasibility unexamined on Apple Silicon & QLoRA and related
PEFT studies {[}6, 7{]} report results exclusively on CUDA-capable GPUs;
the technique's behaviour on Apple's MPS backend, including known
numerical differences in mixed-precision training, is unreported. & All
four adapters were trained to convergence on Apple M4 Pro hardware under
MPS, with the specific adaptations required (bfloat16 precision,
completion-only masking configuration) documented in Chapter 4. \\
Gap 3: No agentic correction loop for UI test scripts & Agentic
self-correction frameworks {[}8, 9, 10{]} have been demonstrated for
general code generation and conversational tasks, but not for the
iterative, static-analysis-guided correction of UI test scripts
specifically. & The BMAD loop was implemented and evaluated across the
complete 279-story dataset for all four fine-tuned adapters, with
acceptance-rate and correction-iteration outcomes reported in Chapter
5. \\
\bottomrule
\end{longtable}
